\section{Chapter 8: Rings}\label{sec:14-appendix-solutions:07-chapter-8:1}

\textbf{1. \texttt{boolAndOrRing} (the ring \(\mathbb{Z}/2\mathbb{Z}\) on \texttt{Bool})}

\begin{lstlisting}
def bool2CommGroup : CommGroup Bool where
  toGroup := boolXorGroup
  comm := by
    intro a b
    cases a with
    | false => cases b with | false => rfl | true => rfl
    | true => cases b with | false => rfl | true => rfl

def bool2Ring : Ring Bool where
  addGrp := bool2CommGroup
  mul := Bool.and
  one := true
  mul_assoc := by
    intro a b c
    cases a with
    | false => cases b with
      | false => cases c with | false => rfl | true => rfl
      | true => cases c with | false => rfl | true => rfl
    | true => cases b with
      | false => cases c with | false => rfl | true => rfl
      | true => cases c with | false => rfl | true => rfl
  one_mul := by
    intro a
    cases a with | false => rfl | true => rfl
  mul_one := by
    intro a
    cases a with | false => rfl | true => rfl
  left_distrib := by
    intro a b c
    cases a with
    | false => cases b with
      | false => cases c with | false => rfl | true => rfl
      | true => cases c with | false => rfl | true => rfl
    | true => cases b with
      | false => cases c with | false => rfl | true => rfl
      | true => cases c with | false => rfl | true => rfl
  right_distrib := by
    intro a b c
    cases a with
    | false => cases b with
      | false => cases c with | false => rfl | true => rfl
      | true => cases c with | false => rfl | true => rfl
    | true => cases b with
      | false => cases c with | false => rfl | true => rfl
      | true => cases c with | false => rfl | true => rfl
\end{lstlisting}

We reuse \texttt{boolXorGroup} from Chapter 6's exercise as the additive group (\texttt{+} = XOR, \texttt{0} = \texttt{false}), and add \texttt{∧} (Boolean and) as multiplication with \texttt{1 = true}. As with \texttt{assoc} in Chapter 6, every axiom is a finite check over \(2\) or \(2^3\) cases, each closed by \href{https://lean-lang.org/doc/reference/latest/Tactic-Proofs/Tactic-Reference/}{\texttt{rfl}} since \texttt{Bool.and}/\texttt{Bool.xor} compute on concrete constructors. This is exactly \(\mathbb{Z}/2\mathbb{Z}\): XOR is addition mod 2, and AND is multiplication mod 2.

\textbf{2. Why both \texttt{left\_\allowbreak{}distrib} and \texttt{right\_\allowbreak{}distrib} are needed}

Distributivity as usually stated, \(a(b+c) = ab+ac\), only lets you expand a product with the \emph{sum on the right}. Its mirror, \((a+b)c = ac+bc\), expands a product with the \emph{sum on the left}. If \texttt{mul} were known to be commutative, these two would be interchangeable --- just apply \texttt{comm} and use the other one. But a general \texttt{Ring} (unlike a \texttt{CommRing}) makes no such assumption. In fact, Chapter 9's proof of \texttt{mul\_\allowbreak{}zero} uses \texttt{left\_\allowbreak{}distrib}, while \texttt{mul\_\allowbreak{}zero\_\allowbreak{}left} (its mirror, an exercise) needs \texttt{right\_\allowbreak{}distrib} instead, precisely because there is no \texttt{mul\_\allowbreak{}comm} field to convert one into the other.

\textbf{3. \texttt{theorem mat2\_\allowbreak{}not\_\allowbreak{}comm : ∃ X Y : Mat2, Mat2.mul X Y ≠ Mat2.mul Y X}}

\begin{lstlisting}
theorem mat2_not_comm : ∃ X Y : Mat2, Mat2.mul X Y ≠ Mat2.mul Y X := by
  refine ⟨X, Y, fun h => ?_⟩
  rw [Mat2.mk.injEq] at h
  exact absurd h.1 (by decide)
\end{lstlisting}

\texttt{X} and \texttt{Y} are the witness pair already computed in the chapter (\texttt{Mat2.mul X Y} evaluates to \texttt{⟨2, 1, 1, 1⟩}, \texttt{Mat2.mul Y X} to \texttt{⟨1, 1, 1, 2⟩}), so \texttt{refine ⟨X, Y, fun h => ?\_\allowbreak{}⟩} only needs to derive \texttt{False} from an assumed equality \texttt{h : Mat2.mul X Y = Mat2.mul Y X}. As the exercise's hint notes, \texttt{by decide} cannot attack \texttt{h} directly --- \texttt{Mat2} has no \texttt{DecidableEq} instance --- so \texttt{Mat2.mk.injEq} first turns \texttt{h} into a conjunction of four \texttt{Int} equalities, one per field. \texttt{h.1}, the first conjunct, states \texttt{2 = 1} (the two matrices' \texttt{a11} entries), which \texttt{Int} \emph{does} have decidable equality for; \texttt{by decide} refutes it directly, and \texttt{absurd} combines the false hypothesis with its refutation to close the goal.
